%!TEX program = xelatex
\documentclass[a4paper,12pt]{article}
\usepackage{fontspec}\defaultfontfeatures{Ligatures=TeX}
\usepackage[a4paper,vmargin={3cm,3.5cm},hmargin={3cm,3cm}]{geometry}
%-----------------------------------------------------------------------------%
\usepackage{settings}
%-----------------------------------------------------------------------------%
%%% Title %%%
\title{G6411 Recitation Notes: \glsentryshort{ols} Finite Sample Properties}
\author{Jesse Chieh Chen\\\href{mailto:cc5229@columbia.edu}{cc5229@columbia.edu}}
\date{\today}
%-----------------------------------------------------------------------------%

\begin{document}

\maketitle

\noindent
Finite sample properties describe an estimator's behavior across repeated samples of a given size.
We ask where its distribution is centered and how dispersed it is.
Unlike purely algebraic properties, these statements require assumptions about how the data are generated.

\section{The Model and Assumptions}

Write the model as
\begin{align}\label{eq-model}
	Y_i=X_i\T\beta+e_i,
	\qquad Y=X\beta+e,
\end{align}
where $X_i\in\reals^k$, $Y=(Y_1,\ldots,Y_n)\T$,
$X=[X_1\;\cdots\;X_n]\T$, and $e=(e_1,\ldots,e_n)\T$.
The coefficient $\beta$ is fixed, while the outcomes, regressors, and errors may be random.
An intercept, if included, counts as one of the $k$ regressors.

\begin{assumption}[Full Rank and Conditional Mean Zero]\label{ass-basic}
	The design matrix $X$ has full column rank $k$, the errors have finite second moments, and
	\begin{align}\label{eq-conditional-mean}
		\E(e\given X)=0.
	\end{align}
\end{assumption}

\begin{remark}
	The conditional mean restriction says that $\E(e_i\given X)=0$ for every observation:
	we condition on the \emph{entire} design matrix, not just $X_i$.
	For \gls{iid} observations, $\E(e_i\given X_i)=0$ gives this restriction.
	Without independent sampling, the two restrictions need not be equivalent.
\end{remark}

Define the conditional error covariance matrix
\begin{align}\label{eq-error-covariance}
	D\coloneqq\E(ee\T\given X)=\var(e\given X).
\end{align}
The equality uses $\E(e\given X)=0$.
Its diagonal entries are $\var(e_i\given X)$, and its off-diagonal entries are $\cov(e_i,e_j\given X)$.

\begin{definition}[Homoskedasticity and Uncorrelated Errors]
	The errors are \emph{conditionally homoskedastic} if all diagonal entries of $D$ equal the same constant $\sigma^2$.
	They are \emph{conditionally uncorrelated} across observations if all off-diagonal entries are zero.
	For time-ordered observations, the latter is also called \emph{no serial correlation}.
	Together, these restrictions give
	\begin{align}\label{eq-spherical-errors}
		D=\sigma^2 I_n.
	\end{align}
\end{definition}

\begin{remark}
	Homoskedasticity concerns the diagonal; absence of correlation concerns the off-diagonal entries.
	Neither condition implies the other, and uncorrelated errors need not be independent.
	We impose \eqref{eq-spherical-errors} only where stated.
\end{remark}

\section{Unbiasedness and Variance}

The starting point for both calculations is
\begin{align}\label{eq-estimation-error}
	\widehat\beta
	=(X\T X)\inv X\T Y
	=\beta+(X\T X)\inv X\T e.
\end{align}
Conditional on $X$, the matrix multiplying $e$ is fixed.

\begin{proposition}[Conditional and Unconditional Unbiasedness]
	Under \autoref{ass-basic},
	\begin{align}
		\E(\widehat\beta\given X)=\beta.
	\end{align}
	Consequently, $\E(\widehat\beta)=\beta$ whenever the unconditional expectation exists.
\end{proposition}
\begin{proof}
	Taking conditional expectations in \eqref{eq-estimation-error} gives
	\begin{align*}
		\E(\widehat\beta\given X)
		=\beta+(X\T X)\inv X\T\E(e\given X)=\beta.
	\end{align*}
	The \gls{lie} then gives $\E(\widehat\beta)=\E[\E(\widehat\beta\given X)]=\beta$.
\end{proof}

\begin{remark}[Orthogonality]
	The population orthogonality restriction $\E(X_i e_i)=0$ does \textsc{not} by itself give finite sample unbiasedness.
	In \eqref{eq-estimation-error},
	$(X\T X)\inv$ is random and generally related to $X\T e$.
	Thus, $\E(X\T e)=0$ does not let us conclude that their product has expectation zero.
	The conditional mean restriction avoids this problem by first holding $X$ fixed.
\end{remark}

\begin{proposition}[Conditional Variance]\label{prop-variance}
	Under \autoref{ass-basic},
	\begin{align}\label{eq-sandwich}
		\var(\widehat\beta\given X)
		=(X\T X)\inv X\T D X(X\T X)\inv.
	\end{align}
	\autoref{eq-sandwich} is also known as the \emph{sandwich formula}.
	If $D=\sigma^2 I_n$, then the sandwich formula simplifies to
	\begin{align}\label{eq-homoskedastic-variance}
		\var(\widehat\beta\given X)=\sigma^2(X\T X)\inv.
	\end{align}
\end{proposition}
\begin{proof}
	\autoref{eq-estimation-error} gives the following:
	\begin{align}
		\var(\widehat\beta\given X)
		&= \var\big((X\T X)\inv X\T e\given X\big)\\
		&= \E\big[(X\T X)\inv X\T e e\T X(X\T X)\inv\given X\big]\\
		&= (X\T X)\inv X\T \E(ee\T\given X) X(X\T X)\inv
		= (X\T X)\inv X\T D X(X\T X)\inv.
	\end{align}
	When $D=\sigma^2 I_n$, the middle cross-product cancels one inverse.
\end{proof}

\begin{remark}[Conditional versus Unconditional Variance]
	When the unconditional second moments exist, the law of total variance gives
	\begin{align*}
		\var(\widehat\beta)
		&=\E[\var(\widehat\beta\given X)]
		+\var[\E(\widehat\beta\given X)]\\
		&=\E[\var(\widehat\beta\given X)].
	\end{align*}
	The second term is zero because the conditional mean is the fixed coefficient $\beta$.
	Under \eqref{eq-spherical-errors}, the unconditional variance is therefore $\sigma^2\E[(X\T X)\inv]$,
	not $\sigma^2[\E(X\T X)]\inv$.
\end{remark}

\section{The Gauss-Markov Theorem}

The Gauss-Markov theorem compares \gls{ols} with other estimators that are linear and unbiased.
We first define the class of estimators and what it means to be best within that class.

\begin{definition}[Best Linear Unbiased Estimator]
	A \emph{linear estimator} has the form $\widetilde\beta=AY$, where $A$ may depend on $X$ but not on $Y$.
	It is \emph{conditionally unbiased} if $\E(\widetilde\beta\given X)=\beta$ for every $\beta$.
	A \emph{best linear unbiased estimator} has the smallest conditional variance for every linear combination of its coefficients among these estimators.
\end{definition}

\begin{remark}[``Best'' and ``Linear'']
	Here
	``best'' means having the smallest conditional variance;
	``linear'' means linear in the outcome vector $Y$, not necessarily in the covariates.
\end{remark}

\begin{remark}[Comparing Covariance Matrices]
	When comparing the variance of two random vectors,
	say $\widehat\beta$ and $\widetilde\beta$,
	it is not sufficient to compare the variance of each component separately.
	Instead, we say that $\widetilde\beta$ has weakly larger variance than $\widehat\beta$ if
	\begin{align}
		\var(\widetilde\beta\given X) - \var(\widehat\beta\given X)
		\succeq 0.
	\end{align}
	That is, the difference is positive semidefinite.
	See \autoref{sec-covariance-ordering}.
\end{remark}

\begin{theorem}[Gauss-Markov]\label{thm-gauss-markov}
	Under \autoref{ass-basic} and $D=\sigma^2 I_n$, \gls{ols} is the best linear unbiased estimator.
\end{theorem}
\begin{proof}
	Unbiasedness for every $\beta$ requires $AX=I_k$ since
	\begin{align*}
		\E(\widetilde\beta\given X)=AX\beta+A\E(e\given X)=AX\beta.
	\end{align*}
	Define $C\coloneqq A-A_0$ where $A_0=(X\T X)\inv X\T$ is the \gls{ols} matrix.
	Since $A_0X=I_k$, we have $CX=0$.
	Consequently, $A_0C\T=(X\T X)\inv(CX)\T=0$.
	Now consider the following:
	\begin{align*}
		\var(\widetilde\beta\given X)
		=\sigma^2(A_0+C)(A_0+C)\T
		=\sigma^2(X\T X)\inv+\sigma^2CC\T.
	\end{align*}
	Note that $\var(\widetilde\beta\given X)-\var(\widehat\beta\given X)=\sigma^2 CC'$
	is positive semidefinite.
\end{proof}

\begin{example}[Unequal Error Variances]
	Suppose $Y_1=\beta+e_1$ and $Y_2=\beta+e_2$, with mean-zero uncorrelated errors of variances $1$ and $4$.
	The \gls{ols} estimate is $(Y_1+Y_2)/2$, with variance $5/4$.
	The estimator $\widetilde\beta=(4Y_1+Y_2)/5$ is also linear and unbiased, but its variance is
	\begin{align*}
		\var(\widetilde\beta)=\frac{16}{25}\cdot1+\frac{1}{25}\cdot4=\frac45<\frac54.
	\end{align*}
	Giving more weight to the less noisy observation improves precision.
	Thus, \gls{ols} need not be best when homoskedasticity fails.
\end{example}

\begin{remark}[What Gauss-Markov Does and Does Not Say]
	The theorem requires neither normality nor independent errors, but it does require $D=\sigma^2 I_n$.
	It compares \gls{ols} only with linear unbiased estimators; it does not rule out improvements from accepting some bias.
	Because unequal variances and correlated errors are common in applications,
	the theorem is merely a benchmark rather than a general efficiency guarantee.
	When \eqref{eq-spherical-errors} fails, the sandwich formula \eqref{eq-sandwich} still applies under \autoref{ass-basic}.
\end{remark}

\section{Estimating the Variance}

Throughout this section, suppose $D=\sigma^2 I_n$ and $n>k$.
Define
\begin{align*}
	P=X(X\T X)\inv X\T,
	\qquad M=I_n-P.
\end{align*}
Since $MX=0$, the residual vector satisfies
\begin{align}\label{eq-residual-error}
	\widehat e=Y-X\widehat\beta=MY=Me.
\end{align}
The residuals are therefore not the original errors: they are the errors after removing their projection onto $X$.

\subsection{Why the Analogue Estimator Is Biased}

If we observed the errors, $n\inv\sum_i e_i^2$ would be an unbiased estimator of $\sigma^2$.
Replacing errors with residuals gives
\begin{align*}
	\widetilde\sigma^2=\frac1n\sum_{i=1}^n\widehat e_i^2=\frac1n e\T M e.
\end{align*}
But fitting the regression removes some of the noise in the sample:
\begin{align*}
	\|e\|^2=\|Pe\|^2+\|Me\|^2.
\end{align*}
Using residuals as if they were errors therefore understates error variance on average.

\begin{proposition}[The Degrees of Freedom Adjustment]\label{prop-error-variance}
	Under \autoref{ass-basic} and $D=\sigma^2 I_n$,
	\begin{align}\label{eq-expected-rss}
		\E(\widehat e\T\widehat e\given X)=(n-k)\sigma^2.
	\end{align}
	Consequently,
	\begin{align}
		\E(\widetilde\sigma^2\given X)=\frac{n-k}{n}\sigma^2,
		\qquad
		\widehat\sigma^2\coloneqq\frac{\widehat e\T\widehat e}{n-k}
	\end{align}
	is conditionally and unconditionally unbiased for $\sigma^2$.
\end{proposition}
\begin{proof}
	For a square matrix $B$, write $\operatorname{tr}(B)$ for the sum of its diagonal entries.
	Using $M\T M=M$ and $\operatorname{tr}(AB)=\operatorname{tr}(BA)$,
	\begin{align*}
		\E(e\T M e\given X)
		=\operatorname{tr}\!\left(M\E(ee\T\given X)\right)
		=\sigma^2\operatorname{tr}(M).
	\end{align*}
	Also,
	\begin{align*}
		\operatorname{tr}(P)
		&=\operatorname{tr}\!\left((X\T X)\inv X\T X\right)=k,\\
		\operatorname{tr}(M)&=n-\operatorname{tr}(P)=n-k.
	\end{align*}
	This proves \eqref{eq-expected-rss}; dividing by $n-k$ gives unbiasedness.
\end{proof}

\begin{remark}[What are Degrees of Freedom?]
	The residual vector satisfies the $k$ independent restrictions $X\T\widehat e=0$.
	It lies in an $(n-k)$-dimensional space: these are its \emph{residual degrees of freedom}.
	The trace calculation counts the same dimensions algebraically.
	Estimating $k$ coefficients leaves $n-k$ directions of noise in the residuals.
\end{remark}

\begin{proposition}[An Unbiased Covariance Estimator]
	Under \autoref{ass-basic} and $D=\sigma^2 I_n$, the matrix
	\begin{align}
		\widehat V=\widehat\sigma^2(X\T X)\inv
	\end{align}
	is conditionally unbiased for $\var(\widehat\beta\given X)$.
\end{proposition}
\begin{proof}
	Conditional on $X$, the inverse matrix is fixed, so
	\begin{align*}
		\E(\widehat V\given X)
		= \E(\widehat\sigma^2\given X)(X\T X)\inv
		= \sigma^2(X\T X)\inv
	\end{align*}
	by \autoref{prop-error-variance}.
\end{proof}

\begin{remark}[Uncorrelated Errors and Correlated Residuals]
	Even when $D=\sigma^2 I_n$,
	the covariance matrix of the residuals need not be diagonal:
	\begin{align}\label{eq-residual-covariance}
		\var(\widehat e\given X)=MDM\T=\sigma^2M.
	\end{align}
	Thus, $\var(\widehat e_i\given X)=\sigma^2(1-P_{ii})$ and, for $i\ne j$,
	$\cov(\widehat e_i,\widehat e_j\given X)=-\sigma^2P_{ij}$.
	The off-diagonal entries need not vanish.
	For an intercept-only regression, $P=\one\one\T/n$, so every distinct pair of residuals has covariance $-\sigma^2/n$.
	Their sum must be zero: a positive residual must be offset by others.
\end{remark}

\begin{remark}[The Role of the Covariance Assumption]
	With a general $D$,
	\begin{align*}
		\E(\widehat e\T\widehat e\given X)=\operatorname{tr}(MD).
	\end{align*}
	Dividing by $n-k$ need not estimate a common error variance, and $\widehat\sigma^2(X\T X)\inv$ need not estimate the sandwich covariance correctly.
	The degrees of freedom correction does not by itself address heteroskedasticity or error correlation.
\end{remark}

\section{Prediction}

For a new observation with regressor vector $X_{n+1}\in\reals^{k\times 1}$, suppose
\begin{align*}
	Y_{n+1}=X_{n+1}\T\beta+e_{n+1},
	\qquad \widehat Y_{n+1}=X_{n+1}\T\widehat\beta.
\end{align*}
Hold $X_{n+1}$ fixed, and suppose the new error is independent of the training sample,
with mean zero and variance $\sigma^2$.
The prediction error decomposes as
\begin{align}\label{eq-prediction-error}
	Y_{n+1}-\widehat Y_{n+1}
	=e_{n+1}-X_{n+1}\T(\widehat\beta-\beta).
\end{align}
The first term is the new observation's genuine error; the second comes from estimating the coefficient.

\begin{proposition}[\Gls{mse} for Prediction]
	Under \autoref{ass-basic}, $D=\sigma^2 I_n$, and the assumptions above,
	for each training observation $i\in\{1,\ldots,n\}$ and the new observation,
	\begin{align}\label{eq-prediction-mse}
		\E[(Y_i-\widehat Y_i)^2\given X] &= \sigma^2-\sigma^2X_i\T(X\T X)\inv X_i,\\
		\E[(Y_{n+1}-\widehat Y_{n+1})^2\given X] &= \sigma^2+\sigma^2X_{n+1}\T(X\T X)\inv X_{n+1}.
	\end{align}
\end{proposition}
\begin{proof}
	Let $\Delta=\widehat\beta-\beta=(X\T X)\inv X\T e$.
	Conditional on $X$, $\E(\Delta\given X)=0$ and
	$\var(\Delta\given X)=\sigma^2(X\T X)\inv$.
	For a training observation, $Y_i-\widehat Y_i=e_i-X_i\T\Delta$.
	Since $\cov(e,e_i\given X)$ has entry $\sigma^2$ at position $i$ and zero elsewhere,
	\begin{align*}
		\cov(e_i,X_i\T\Delta\given X)
		=\sigma^2X_i\T(X\T X)\inv X_i.
	\end{align*}
	Thus,
	\begin{align*}
		\E[(Y_i-\widehat Y_i)^2\given X]
		&=\var(e_i\given X)+\var(X_i\T\Delta\given X)\\
		&\quad-2\cov(e_i,X_i\T\Delta\given X)\\
		&=\sigma^2+\sigma^2X_i\T(X\T X)\inv X_i
		-2\sigma^2X_i\T(X\T X)\inv X_i\\
		&=\sigma^2-\sigma^2X_i\T(X\T X)\inv X_i.
	\end{align*}
	For the new observation, $e_{n+1}$ is independent of the training sample, so the covariance term is zero.
	Both terms in \eqref{eq-prediction-error} have conditional mean zero, giving
	\begin{align*}
		\E[(Y_{n+1}-\widehat Y_{n+1})^2\given X]
		&=\var(e_{n+1}\given X)+\var(X_{n+1}\T\Delta\given X)\\
		&=\sigma^2+\sigma^2X_{n+1}\T(X\T X)\inv X_{n+1}.
	\end{align*}
\end{proof}

\begin{remark}[Predicting an Outcome versus Its Mean]
	For the conditional mean $X_{n+1}\T\beta$, the estimation \gls{mse} is just
	$\sigma^2X_{n+1}\T(X\T X)\inv X_{n+1}$.
	Predicting an individual outcome adds the irreducible error variance $\sigma^2$.
	At a training regressor value $X_{n+1}=X_i$, the new-outcome prediction error has expected square $\sigma^2(1+P_{ii})$,
	whereas the training residual has expected square $\sigma^2(1-P_{ii})$.
	The training residual benefits from having used $Y_i$ to estimate the coefficient.
\end{remark}

\section{Omitted and Unnecessary Variables}

Partition the design matrix as $X=[X_1\;X_2]$ with full column rank.
Here $X_1$ and $X_2$ are blocks of columns, not individual observations.
For the variance comparisons in this section, assume that the true errors satisfy
$\E(e\given X)=0$ and $\var(e\given X)=\sigma^2 I_n$.
We compare estimates on the same sample, conditional on the full $X$.
Define
\begin{align*}
	M_2=I_n-X_2(X_2\T X_2)\inv X_2\T.
\end{align*}

\subsection{Omitting Relevant Variables}

Suppose the true model is
\begin{align*}
	Y=X_1\beta_1+X_2\beta_2+e,
\end{align*}
but we regress $Y$ only on $X_1$.

\paragraph{Bias.}

The omitted contribution becomes part of the regression error $u=X_2\beta_2+e$.
The short-regression estimate satisfies
\begin{align}\label{eq-omitted-variable-bias}
	\widetilde\beta_1
	=(X_1\T X_1)\inv X_1\T Y
	=\beta_1+(X_1\T X_1)\inv X_1\T X_2\beta_2
	+(X_1\T X_1)\inv X_1\T e.
\end{align}
Its conditional bias is therefore
\begin{align}
	\E(\widetilde\beta_1\given X)-\beta_1
	=(X_1\T X_1)\inv X_1\T X_2\beta_2.
\end{align}
Bias arises when the omitted contribution $X_2\beta_2$ projects onto the included regressors.
It disappears if $\beta_2=0$ or $X_1\T X_2=0$, though other cancellations are also possible.

\begin{example}[The Usual Omitted Variable Bias Formula]
	Let the true model be $Y_i=\alpha+\beta x_i+\gamma z_i+e_i$,
	and regress $Y_i$ on an intercept and $x_i$ alone.
	Conditional on the observed $x_i$ and $z_i$, the slope's bias is
	\begin{align*}
		\gamma\frac{\sum_i(x_i-\bar x)(z_i-\bar z)}{\sum_i(x_i-\bar x)^2}.
	\end{align*}
	For there to be omitted variable bias it must be the case that
	(1) $z_i$ affects the outcome, i.e., $\gamma\neq0$ and
	(2) it is related to the included regressor, i.e., $\sum_{i=1}^{n}(x_i-\bar x)(z_i-\bar z)\neq0$.
\end{example}

\paragraph{Variance.}

Conditional on the full $X$, the omitted contribution is fixed, so
\begin{align}
	\var(\widetilde\beta_1\given X)&=\sigma^2(X_1\T X_1)\inv,\\
	\var(\widehat\beta_1\given X)&=\sigma^2(X_1\T M_2X_1)\inv,
\end{align}
where $\widehat\beta_1$ comes from the full regression.
The second identity follows by the \gls{fwl} theorem.
Since $X_1\T M_2X_1\preceq X_1\T X_1$, inversion reverses the ordering:
\begin{align}\label{eq-variable-variance-comparison}
	\var(\widetilde\beta_1\given X)
	\preceq\var(\widehat\beta_1\given X).
\end{align}
\autoref{sec-covariance-ordering} explains this covariance ordering.
Omitting variables can reduce variance while introducing bias.
This comparison uses the same true $\sigma^2$ in both models;
it does not say that the short regression's reported standard errors correctly measure uncertainty about $\beta_1$.

\begin{remark}[When are the Two Equal?]
	For $\sigma^2>0$, the variance comparison \eqref{eq-variable-variance-comparison} is an equality if and only if $X_1\T M_2X_1=X_1\T X_1$.
	To see when this occurs, note that
	\begin{align*}
		X_1\T X_1-X_1\T M_2X_1
		=(X_2\T X_1)\T(X_2\T X_2)\inv(X_2\T X_1).
	\end{align*}
	Since $(X_2\T X_2)\inv$ is positive definite, this is zero exactly when $X_1\T X_2=0$.
	Thus, equality holds when the two blocks are orthogonal:
	controlling for $X_2$ removes no variation from $X_1$.
	Here orthogonality refers to sample cross-products; it is not the same as zero population correlation.
\end{remark}

\subsection{Including Unnecessary Variables}

Now suppose the true model is $Y=X_1\beta_1+e$, but we also include $X_2$.
The expanded model is correctly specified with $\beta_2=0$.

\paragraph{Bias.}

Both the short and full estimates of $\beta_1$ are conditionally unbiased
because both models are correctly specified and $\E(e\given X)=0$.

\paragraph{Variance.}

If $X_2$ and $X_1$ are not orthogonal,
then the inclusion of $X_2$ ``soaks up'' some usable variation in $X_1$,
i.e., the relation \eqref{eq-variable-variance-comparison} still holds.
This means that by including unnecessary variables, we can increase the variance of the estimate of $\beta_1$.

\subsection{The Bias-Variance Tradeoff}

To compare a biased estimate with an unbiased one, we need a measure that accounts for both bias and dispersion.
\Gls{mse} measures the expected squared distance from the parameter.

\begin{proposition}[\Gls{mse} Decomposition]
	For a scalar estimator $\widetilde\theta$ of a fixed parameter $\theta$, we have
	\begin{align}
		\E[(\widetilde\theta-\theta)^2\given X]
		=\var(\widetilde\theta\given X)
		+\left(\E(\widetilde\theta\given X)-\theta\right)^2.
	\end{align}
\end{proposition}
\begin{proof}
	Consider the following:
	\begin{align}
		\widetilde\theta-\theta
		= (\widetilde\theta-\E(\widetilde\theta\given X))+(\E(\widetilde\theta\given X)-\theta).
	\end{align}
	Square both sides and take conditional expectations and we have
	\begin{align*}
		\E[(\widetilde\theta-\theta)^2\given X]
		&= \E[(\widetilde\theta-\E(\widetilde\theta\given X))^2\given X]
		+ (\E(\widetilde\theta\given X)-\theta)^2\\
		&\quad+ 2\underbrace{\E[(\widetilde\theta-\E(\widetilde\theta\given X))(\E(\widetilde\theta\given X)-\theta)\given X]}_{=0} \\
		&= \var(\widetilde\theta\given X)
		+ (\E(\widetilde\theta\given X)-\theta)^2.
		\qedhere
	\end{align*}
\end{proof}

\begin{remark}[Vector Estimators]
	For a vector estimator we can apply the same decomposition component-by-component.
	For example, we have
	\begin{align*}
		\E[\|\widetilde\beta-\beta\|^2\given X]
		=\operatorname{tr}[\var(\widetilde\beta\given X)]+\|\text{bias}\|^2
	\end{align*}
	where $\text{bias}=\E(\widetilde\beta\given X)-\beta$.
\end{remark}

\begin{remark}[The Bias-Variance Tradeoff]
	A biased estimator can have smaller \gls{mse} if the variance reduction outweighs squared bias.
	Neither including every available variable nor choosing the smallest model is always best.
	This motivates us to define \emph{regularization} methods that explore this tradeoff.
\end{remark}

\section{Ridge Regression}

Ridge regression adds a penalty for large coefficients to the least squares objective.
This stabilizes estimation by shrinking the coefficient directions that \gls{ols} estimates least precisely.
The multicollinearity and singular value interpretation is developed in \autoref{sec-svd}.

\begin{definition}[Ridge Regression]
	For $\lambda>0$, the ridge estimator minimizes
	\begin{align}\label{eq-ridge-loss}
		\widehat\beta_\lambda
		=\arg\min_{b\in\reals^k}
		\left\{\|Y-Xb\|^2+\lambda\|b\|^2\right\}.
	\end{align}
	Its solution is
	\begin{align}\label{eq-ridge-solution}
		\widehat\beta_\lambda=(X\T X+\lambda I_k)\inv X\T Y.
	\end{align}
\end{definition}
\begin{proof}
	The first order condition is $(X\T X+\lambda I_k)b=X\T Y$.
	For $\lambda>0$, the coefficient matrix is positive definite even if $X$ is rank deficient,
	so the minimizer is unique.
\end{proof}

\begin{remark}[The Intercept and Units]
	Equation~\eqref{eq-ridge-loss} penalizes every coefficient, including an intercept if present.
	Usually we leave the intercept unpenalized, center the outcome and regressors, and apply ridge to the slopes.
	The penalty also depends on regressor units, so scaling the regressors is part of specifying the estimator.
	The calculations below use the penalty in \eqref{eq-ridge-loss} and a fixed $\lambda$.
\end{remark}

\begin{proposition}[Ridge Bias and Variance]\label{prop-ridge}
	Under \autoref{ass-basic} and $D=\sigma^2 I_n$, for fixed $\lambda>0$,
	write $Q=X\T X$ and $B_\lambda=(Q+\lambda I_k)\inv$.
	\begin{align}
		\E(\widehat\beta_\lambda\given X)-\beta
		&=-\lambda B_\lambda\beta,\\
		\var(\widehat\beta_\lambda\given X)
		&=\sigma^2B_\lambda Q B_\lambda.
	\end{align}
	In particular,
	$\var(\widehat\beta_\lambda\given X)\preceq\var(\widehat\beta\given X)$.
\end{proposition}
\begin{proof}
	Substitute $Y=X\beta+e$ into the ridge formula:
	\begin{align*}
		\widehat\beta_\lambda
		=B_\lambda X\T(X\beta+e)
		=B_\lambda Q\beta+B_\lambda X\T e.
	\end{align*}
	Since $B_\lambda\inv=Q+\lambda I_k$, we have
	\begin{align*}
		B_\lambda Q
		=B_\lambda(B_\lambda\inv-\lambda I_k)
		=I_k-\lambda B_\lambda.
	\end{align*}
	Using $\E(e\given X)=0$, the conditional bias is therefore
	\begin{align*}
		\E(\widehat\beta_\lambda\given X)-\beta
		=(B_\lambda Q-I_k)\beta+B_\lambda X\T\E(e\given X)
		=-\lambda B_\lambda\beta.
	\end{align*}
	For the variance, $B_\lambda Q\beta$ is fixed conditional on $X$.
	Also, $B_\lambda$ is symmetric because $Q+\lambda I_k$ is symmetric.
	Thus,
	\begin{align*}
		\var(\widehat\beta_\lambda\given X)
		&=\var(B_\lambda X\T e\given X)\\
		&=B_\lambda X\T\var(e\given X)X B_\lambda\T
		=\sigma^2B_\lambda X\T X B_\lambda
		=\sigma^2B_\lambda Q B_\lambda.
	\end{align*}
	To compare with \gls{ols}, insert $I_k=B_\lambda B_\lambda\inv=B_\lambda\inv B_\lambda$ around $Q\inv$ and factor out $B_\lambda$:
	\begin{align*}
		\var(\widehat\beta\given X)-\var(\widehat\beta_\lambda\given X)
		&=\sigma^2(Q\inv-B_\lambda Q B_\lambda)\\
		&=\sigma^2B_\lambda\big[(Q+\lambda I_k)Q\inv(Q+\lambda I_k)-Q\big]B_\lambda\\
		&=\sigma^2B_\lambda\big[Q+2\lambda I_k+\lambda^2Q\inv-Q\big]B_\lambda\\
		&=\sigma^2B_\lambda(2\lambda I_k+\lambda^2Q\inv)B_\lambda.
	\end{align*}
	Since $Q\inv$ is positive definite and $\lambda>0$, the middle matrix $2\lambda I_k+\lambda^2Q\inv$ is positive definite.
	Multiplying it on the left by $B_\lambda$ and on the right by $B_\lambda\T=B_\lambda$ preserves positive semidefiniteness,
	proving the covariance ordering.
\end{proof}

\begin{remark}[Does Ridge Improve \gls{mse}?]
	With the notation in \autoref{prop-ridge}, the conditional \glspl{mse} are
	\begin{align*}
		R_{\mathrm{OLS}}\coloneqq\E[\|\widehat\beta-\beta\|^2\given X]
		&=\sigma^2\operatorname{tr}(Q\inv),\\
		R_\lambda\coloneqq\E[\|\widehat\beta_\lambda-\beta\|^2\given X]
		&=\sigma^2\operatorname{tr}(B_\lambda Q B_\lambda)
		+\lambda^2\|B_\lambda\beta\|^2.
	\end{align*}
	Their difference is
	\begin{align*}
		R_\lambda-R_{\mathrm{OLS}}
		=\underbrace{\lambda^2\|B_\lambda\beta\|^2}_{\text{squared bias}}
		-\underbrace{\lambda\sigma^2\operatorname{tr}\!\left[B_\lambda(2I_k+\lambda Q\inv)B_\lambda\right]}_{\text{variance reduction}}.
	\end{align*}
	Ridge improves \gls{mse} exactly when variance reduction exceeds squared bias.
	For $\sigma^2>0$, $\lambda\|\beta\|^2<2\sigma^2$ is sufficient:
	using $\|B_\lambda\beta\|^2\le\|\beta\|^2\operatorname{tr}(B_\lambda^2)$ gives
	\begin{align*}
		R_\lambda-R_{\mathrm{OLS}}
		\le\lambda\big(\lambda\|\beta\|^2-2\sigma^2\big)\operatorname{tr}(B_\lambda^2)<0.
	\end{align*}
	Thus, a sufficiently small positive penalty improves coefficient \gls{mse} for any fixed $X$ and $\beta$.
	The bound depends on unknown parameters; not every penalty improves \gls{mse}.
	Ridge is generally biased, so this is consistent with \nameref{thm-gauss-markov}.
\end{remark}

\section{From Moments to Finite Sample Inference}

Knowing the mean and variance of $\widehat\beta$ does not determine its full distribution.
For example, an additional normality assumption gives the exact conditional law:
\begin{align}
	e\given X&\sim\normaldistribution(0,\sigma^2 I_n),\label{eq-normal-errors}\\
	\widehat\beta\given X&\sim\normaldistribution\!\left(\beta,\sigma^2(X\T X)\inv\right).\notag
\end{align}
The usual exact normal and $t$ inference results rely on this additional distributional assumption.
Large sample theory instead derives an approximate distribution for the estimator as the sample size grows.
Under suitable sampling and moment conditions, this permits approximate inference without requiring exactly normal errors.

\vfill
\section*{Acronyms}
\printglossaries

\clearpage
\appendix
\section{How to Compare Covariance Matrices}\label{sec-covariance-ordering}

\begin{definition}[Positive Semidefinite Ordering]
	A symmetric matrix $B$ is \emph{positive semidefinite} if $a\T Ba\ge0$ for every vector $a$.
	For symmetric matrices $V_1$ and $V_2$ of the same size, write
	\begin{align*}
		V_1\preceq V_2
	\end{align*}
	if $V_2-V_1$ is positive semidefinite.
\end{definition}

\begin{remark}
	For covariance matrices, this means that every linear combination has at least as much variance under $V_2$:
	\begin{align*}
		a\T V_1a\le a\T V_2a\quad\text{for every }a.
	\end{align*}
	In particular, every diagonal variance is larger or equal, but that condition alone does not establish the ordering.
	Any matrix $CC\T$ is positive semidefinite because $a\T CC\T a=\|C\T a\|^2$.
	This is the comparison used in the Gauss-Markov proof.
\end{remark}

\begin{proposition}[Inversion Reverses the Ordering]
	If $B$ and $C$ are positive definite and $B\preceq C$, then $C\inv\preceq B\inv$.
\end{proposition}
\begin{proof}
	Let $H=B^{-1/2}CB^{-1/2}$, where the symmetric square root is obtained by taking square roots of the eigenvalues.
	Then $H\succeq I$, so its eigenvalues are at least one and $H\inv\preceq I$.
	Therefore,
	\begin{align*}
		C\inv=B^{-1/2}H\inv B^{-1/2}\preceq B\inv.
	\end{align*}
	The last step uses the fact that multiplying a positive semidefinite matrix on the left by $R$ and on the right by $R\T$ preserves positive semidefiniteness.
\end{proof}

In the variable-selection comparison, $I_n-M_2$ is a projection matrix, so
\begin{align*}
	X_1\T X_1-X_1\T M_2X_1
	=X_1\T(I_n-M_2)X_1\succeq0.
\end{align*}
Both cross-product matrices are positive definite under full column rank of $[X_1\;X_2]$.
The inverse ordering therefore gives \eqref{eq-variable-variance-comparison}.

\section{Multicollinearity and the Singular Value Decomposition}\label{sec-svd}

\subsection{What Makes \glsentryshort{ols} Unstable?}

Exact multicollinearity means that $X$ lacks full column rank.
Near multicollinearity means that, after appropriate scaling, some columns are nearly linear combinations of the others.
The coefficients may then be very imprecisely estimated even when the fitted values are stable.

For full-column-rank $X$, write its full singular value decomposition as
\begin{align}\label{eq-svd}
	X=USV\T,
	\qquad S=\begin{bmatrix}
		\diag(s_1,\ldots,s_k)\\
		0_{(n-k)\times k}
	\end{bmatrix},
\end{align}
where $U\in\reals^{n\times n}$, $S\in\reals^{n\times k}$, and $V\in\reals^{k\times k}$,
with $s_1\ge\cdots\ge s_k>0$.
Both $U$ and $V$ are orthogonal matrices:
$U\T U=UU\T=I_n$ and $V\T V=VV\T=I_k$.
We use $S$ rather than $D$ to distinguish singular values from the error covariance matrix.

\begin{remark}[The Left Singular Vectors]
	The columns $u_1,\ldots,u_n$ of $U$ form an orthonormal basis for $\reals^n$.
	The first $k$ span the column space of $X$.
	The remaining $n-k$ are orthogonal to that space and do not contribute to $X$, because the corresponding rows of $S$ are zero.
\end{remark}

Since $S$ is rectangular, we use $S\T S$, not $S^2$:
\begin{align*}
	S\T S&=\diag(s_1^2,\ldots,s_k^2),\\
	X\T X&=V(S\T S)V\T,\\
	(X\T X)\inv&=V(S\T S)\inv V\T,\\
	\widehat\beta&=V(S\T S)\inv S\T U\T Y.
\end{align*}
The columns $v_j$ of $V$ are eigenvectors of $X\T X$, with eigenvalues $s_j^2$.
For $j=1,\ldots,k$, the corresponding column $u_j$ of $U$ gives
\begin{align}\label{eq-svd-coefficients}
	v_j\T\widehat\beta=\frac{u_j\T Y}{s_j},
	\qquad
	\var(v_j\T\widehat\beta\given X)=\frac{\sigma^2}{s_j^2}
\end{align}
under $D=\sigma^2 I_n$.
Small singular values amplify the noise when recovering coefficients.

\begin{remark}[Condition Indices and Scaling]
	The condition index for direction $j$ is $s_1/s_j$;
	the largest is the condition number $\kappa(X)=s_1/s_k$.
	Equivalently, it is the square root of the ratio of the largest and smallest eigenvalues of $X\T X$.
	A large value indicates a poorly determined coefficient direction.
	These quantities depend on units, so compare suitably scaled regressor columns rather than treating different measurement units as evidence of collinearity.
\end{remark}

\begin{remark}[Projection versus Recovering Coefficients]
	The projection keeps the first $k$ left-singular-vector directions and removes the remaining $n-k$:
	\begin{align*}
		P=U\begin{bmatrix}I_k&0\\0&0_{(n-k)\times(n-k)}\end{bmatrix}U\T.
	\end{align*}
	Thus, the fitted vector and residual vector are
	\begin{align*}
		\widehat Y&=PY=\sum_{j=1}^k u_j(u_j\T Y),\\
		\widehat e&=MY=\sum_{j=k+1}^n u_j(u_j\T Y).
	\end{align*}
	Thus, \gls{ols} keeps the components of $Y$ along the orthonormal directions spanning the columns of $X$.
	Recovering coefficients requires an additional division by $s_j$ and rotation by $V$.
	It is this division, not the projection itself, that becomes unstable when $s_j$ is small.
\end{remark}

\subsection{Ridge in Singular-Value Directions}

The singular value decomposition gives
\begin{align}\label{eq-ridge-svd}
	\widehat\beta_\lambda
	&=V(S\T S+\lambda I_k)\inv S\T U\T Y\\
	&=\sum_{j=1}^k v_j\frac{s_j}{s_j^2+\lambda}(u_j\T Y).
\end{align}
Compared with \eqref{eq-svd-coefficients}, ridge multiplies the \gls{ols} coefficient in direction $v_j$ by
\begin{align*}
	\frac{s_j^2}{s_j^2+\lambda}.
\end{align*}
The smallest singular-value directions receive the strongest shrinkage.
For the fitted values,
\begin{align*}
	\widehat Y_\lambda
	=\sum_{j=1}^k u_j\frac{s_j^2}{s_j^2+\lambda}(u_j\T Y).
\end{align*}
Unlike the \gls{ols} projection, ridge retains only a fraction of each outcome component in the regressor space.

\begin{remark}[Variance and \gls{mse} by Direction]
	Under \autoref{ass-basic} and $D=\sigma^2 I_n$, write $\theta_j=v_j\T\beta$.
	For fixed $\lambda>0$,
	\begin{align*}
		\E(v_j\T\widehat\beta_\lambda\given X)-\theta_j
		&=-\frac{\lambda\theta_j}{s_j^2+\lambda},\\
		\var(v_j\T\widehat\beta_\lambda\given X)
		&=\frac{\sigma^2s_j^2}{(s_j^2+\lambda)^2}
		\le\frac{\sigma^2}{s_j^2},\\
		\E[(v_j\T\widehat\beta_\lambda-\theta_j)^2\given X]
		&=\frac{\lambda^2\theta_j^2+\sigma^2s_j^2}{(s_j^2+\lambda)^2}.
	\end{align*}
	These formulas show the variance reduction and squared bias separately in each direction.
\end{remark}

\end{document}
